\section{Cross-paper synthesis}
\label{sec:synthesis}

\subsection{The explanatory chain}

The literature does not support a single statement such as ``HBM makes
LLMs faster''.  It supports a chain of conditional propositions.  A
workload phase determines the data reused and the operations issued;
the live set and placement determine which memory tier serves those
accesses; and the resulting bandwidth, latency and matrix throughput
determine a user-visible metric.  Skipping any link turns a measured
correlation into an unsupported cause.

\begin{figure}[htbp]
\centering
\begin{tikzpicture}[
  node distance=5mm,
  >=Latex,
  chainbox/.style={
    draw=black!55,
    fill=neutralcol,
    rounded corners=2pt,
    text width=2.45cm,
    minimum height=1.25cm,
    align=center,
    font=\small
  },
  hbmstep/.style={chainbox, fill=hbmcol!15, draw=hbmcol!80!black},
  metric/.style={chainbox, fill=accent!8, draw=accent}
]
  \node[chainbox] (phase) {workload phase\\prefill or decode};
  \node[chainbox, right=of phase] (set) {verified live set\\weights, KV, scratch};
  \node[hbmstep, right=of set] (place) {residency\\local HBM or DDR};
  \node[chainbox, right=of place] (resource) {achieved resource\\bandwidth or compute};
  \node[metric, right=of resource] (metric) {reported metric\\TTFT, ITL, tokens/s};
  \draw[->, thick] (phase) -- (set);
  \draw[->, thick] (set) -- (place);
  \draw[->, thick] (place) -- (resource);
  \draw[->, thick] (resource) -- (metric);
\end{tikzpicture}
\caption{Causal chain required for an interpretable memory-tier result.
The experiment must observe, rather than assume, the central three
links.}
\label{fig:causal-chain}
\end{figure}

For a phase with work \(W\), data movement \(Q\), peak relevant compute
\(P\), and sustained local bandwidth \(B\), a Roofline-style lower bound
is~\cite{williams2009roofline}
\[
  T_{\mathrm{phase}} \;\ge\;
  \max\!\left(\frac{W}{P},\,\frac{Q}{B}\right)
\]
The equation is a diagnostic model, not
a prediction of the HBM speedup.  \(Q\) depends on quantization, cache
reuse, batching, context, packing and runtime behaviour; \(B\) must be
measured under the same placement and core binding as the inference
run.

\subsection{Conclusions supported by more than one source}

\paragraph{HBM is a conditional resource.}
Ibeid et al.\ show a large sustained-bandwidth advantage for local HBM,
but also a DDR latency advantage at low load and capacity-dependent
application behaviour.  Na et al.\ show that an LLM runtime can lose
performance when it fails to respect topology.  Goto et al.\ use DDR
for the oversized global matrix and HBM only for staged CPU-side panel
data.  Together, these results support deliberate placement based on
phase and live-set size.  They do not support an all-HBM policy for
every allocation.

\paragraph{Prefill and decode must be reported separately.}
Prefill processes many prompt tokens and exposes comparatively large
matrix operations; autoregressive decode repeatedly streams weights for
few newly generated tokens and grows the KV state with context and
batch.  Their relative pressure on AMX and memory bandwidth can differ.
One end-to-end average can therefore hide the mechanism.  Time to first
token (TTFT), inter-token latency (ITL) and phase throughput are separate
outcomes.

\paragraph{Capacity is a live-set property.}
Nominal model-file size is not sufficient.  Runtime-packed weights, KV
state, graph and scratch buffers, allocator overhead and page cache can
change resident memory.  Goto et al.'s explicit DDR/HBM staging makes
the general issue concrete: the bulk object can exceed the high-bandwidth
tier even when a critical subproblem fits.  The thesis should measure
residency and peak allocation before labelling a run ``within 64\,GB''.

\paragraph{Topology is part of the treatment.}
The papers' apparently different SNC4 outcomes are not a hardware
contradiction.  Ibeid et al.\ obtain high SNC4 bandwidth with explicit
domain placement and recover MPI performance through network mapping.
Na et al.\ observe remote traffic and lower inference performance with
the runtime setup they test.  Goto et al.\ co-locate ranks, cores,
memory, GPUs and NICs to prevent cross-socket traffic.  SNC4 exposes
locality; it does not make software local automatically.

\paragraph{CPU--GPU comparisons are configuration-specific.}
Na et al.\ compare a direct CPU path with single-GPU FlexGen offload.
Goto et al.\ instead demonstrate a highly tuned, heterogeneous CPU--GPU
pipeline.  These are different questions.  Neither supports a general
claim that one processor class is superior.  A GPU comparison is not
needed to answer the primary HBM-versus-DDR question.

\subsection{What the Aurora deployment paper adds}

Goto et al.\ supply a detailed production example of explicit tier
staging.  They also strengthen three methodological requirements:

\begin{enumerate}
  \item report the exact mapping of compute, memory and communication
        resources;
  \item analyse phase time, because the dominant bottleneck changes as
        execution progresses; and
  \item distinguish observed campaign outcomes from controlled
        ablations.
\end{enumerate}

The same paper also sets a limit on the analogy.  HPL's bulk DDR matrix,
HBM-staged panels, GPU GEMM, MPI communication and global scale differ
from a single-node \lcpp{} decode loop.  The paper suggests that
selective tiering is worth testing; it does not establish which LLM
tensors should occupy HBM or how much improvement to expect.

\subsection{Claims that should not enter the thesis as facts}

The following statements are either too broad or not identified by the
reviewed evidence:

\begin{itemize}
  \item that HBM makes inference faster by the STREAM bandwidth ratio;
  \item that one socket is inherently better than two sockets;
  \item that SNC4 is inherently better or worse than Quadrant mode;
  \item that CPUs outperform GPUs for LLM inference in general;
  \item that a mixed-precision HPL/HPL-MxP ratio measures AMX speedup;
  \item that \lcpp{} is already the cause of a performance limitation
        before a correctly placed baseline is run; or
  \item that a gap in the supplied corpus is a field-wide novelty claim.
\end{itemize}

\subsection{From evidence to candidate questions}

Table~\ref{tab:evidence-questions} separates the direct extension of
existing evidence from more speculative work.

\begin{table}[H]
\centering
\small
\caption{How the reviewed evidence motivates, but does not answer, the
candidate thesis questions.}
\label{tab:evidence-questions}
\begin{tabularx}{\textwidth}{@{}P{3.1cm}YY@{}}
\toprule
\textbf{Candidate question} & \textbf{Existing basis} &
\textbf{Missing observation} \\
\midrule
Same-socket HBM versus DDR for quantized inference
& Xeon Max memory bandwidth and BF16 inference are each measured
  separately.
& Causal tier comparison with identical socket, runtime, model,
  quantization, binding and phase. \\
\addlinespace
Capacity and KV boundary
& HBM capacity effects occur in HPC; LLM KV state grows with context
  and batch.
& Measured live-set crossing with weights and KV residency identified
  separately. \\
\addlinespace
Topology-aware scaling
& Explicit SNC4 placement helps HPC; unmanaged locality harms the
  reported LLM run.
& Controlled local, remote, replica and partitioned designs under one
  software version. \\
\addlinespace
Independent memory-tier placement
& Aurora HPL stages selected data between DDR and HBM.
& A CPU LLM runtime control and an ablation showing that selective
  placement beats the best whole-process policy. \\
\addlinespace
Energy per token
& The position paper motivates efficiency.
& A validated measurement boundary and uncertainty for node or socket
  energy. \\
\bottomrule
\end{tabularx}
\end{table}

The first question is the strongest because it follows directly from
the corpus and can be answered without modifying the runtime.  The
others become valuable only after the baseline identifies where the
working set, remote traffic or software control is limiting the result.
